\documentclass[11pt]{article}

\usepackage[margin=1in]{geometry}
\usepackage{amsmath,amssymb,amsthm}
\usepackage{booktabs}
\usepackage{hyperref}
\usepackage{xcolor}
\usepackage{enumitem}

\hypersetup{colorlinks=true, linkcolor=blue!60!black, citecolor=blue!60!black, urlcolor=blue!60!black}

\title{\textbf{Can Language Models Learn a PDE?}\\[4pt]
\large Probing Scientific Reasoning via Weak-Form Identification\\of an Unknown Elliptic Equation}
\author{}
\date{}

\newcommand{\dd}{\,\mathrm{d}}

\begin{document}
\maketitle

\begin{abstract}
We present a benchmark that tests whether a large language model can
identify an unknown partial differential equation from data, using the same
methodology a mathematician would: probing the equation with test functions
through its weak formulation, analyzing the resulting integrals, and
recovering the unknown coefficients.  The model receives solution data, a
budget of queries that return weak-form integrals, and free access to
diagnostic tools---but must design its own experimental strategy.  We show
that a single paragraph of mathematical context about ill-conditioning,
added to the initial prompt, reduces coefficient recovery error by 42\% on a
hard problem instance, with the model demonstrating qualitatively different
reasoning: checking weight balances before committing queries, seeking test
functions that emphasize underconstrained terms, and articulating why certain
measurements are more informative than others.  The benchmark cleanly
separates tool use from scientific reasoning, providing a measurable test of
whether language models can \emph{do science}---form hypotheses, design
experiments, and diagnose failures---rather than merely follow instructions.
\end{abstract}

%----------------------------------------------------------------------
\section{Introduction}
%----------------------------------------------------------------------

Benchmarks for scientific reasoning in large language models (LLMs) often
test whether the model can \emph{recall} known facts or \emph{solve}
textbook problems.  These tasks measure knowledge retrieval, not the open-ended
reasoning that characterizes real scientific work: forming hypotheses,
designing experiments to test them, interpreting results, and iterating.

We propose a benchmark built around a classical problem in applied
mathematics: \emph{identifying an unknown elliptic PDE from its solution}.
The model is given the solution $u(x)$ on a grid, a budget of queries that
probe the PDE through its weak formulation, and a set of computational
tools.  It must design a sequence of test functions, interpret the
resulting integrals, and recover the unknown coefficient functions.

The key insight is that this problem is \emph{inherently ill-conditioned}
in a way that rewards scientific reasoning.  Not all test functions are
equally informative: some primarily constrain the diffusion coefficient
while revealing almost nothing about the reaction coefficient.  A model
that mechanically collects equations and solves the linear system will get
a mediocre answer.  A model that reasons about \emph{which measurements
are informative}---checking weight balances, seeking test functions that
emphasize underconstrained terms---will do significantly better.

Our experimental design isolates this reasoning capability.  We compare
two runs that differ \textbf{only in the initial prompt}: a baseline
(``Begin.'') and an enhanced version that includes one paragraph explaining
\emph{why} the weak-form system can be ill-conditioned.  The tools, the
PDE, the query budget, and all code are identical.  The enhanced prompt
produces a 42\% reduction in coefficient error, with the model exhibiting
qualitatively different experimental design.

%----------------------------------------------------------------------
\section{Problem Setup}
%----------------------------------------------------------------------

\subsection{The Unknown PDE}

The model is asked to identify the coefficient functions $a(x)$, $b(x)$,
$c(x)$, and $f(x)$ in the second-order elliptic boundary value problem
\begin{equation}\label{eq:pde}
    -\bigl(a(x)\, u'(x)\bigr)' + b(x)\, u'(x) + c(x)\, u(x) = f(x),
    \qquad x \in [0,1],
\end{equation}
with homogeneous Dirichlet conditions $u(0) = u(1) = 0$.  Each coefficient
is a linear combination of shifted Legendre polynomials:
\[
    a(x) = \sum_{j=0}^{N-1} a_j\, P_j(x),
    \qquad
    b(x) = \sum_{j=0}^{N-1} b_j\, P_j(x),
    \qquad \text{etc.}
\]
The model is told that each coefficient is ``smooth and represented by $N$
parameters,'' giving $4N$ unknowns total.  For the \emph{hard} difficulty
used in our experiments, $N = 4$ and there are 16 unknowns.  The
coefficients are drawn randomly with controlled sparsity and scale.

\subsection{What the Model Receives}

The model is given the following information and tools:

\paragraph{Solution data.}
The numerical solution $u(x)$ evaluated on a grid of points.  The model
can also request $u(x)$ and $u'(x)$ at arbitrary points via a free
\texttt{COMPUTE: eval\_solution} command.

\paragraph{The weak form identity.}
The system prompt states the weak formulation obtained by integration by
parts: for any test function $\varphi$ with $\varphi(0) = \varphi(1) = 0$,
\begin{equation}\label{eq:weak}
    \int_0^1 a(x)\, u'(x)\, \varphi'(x)\dd x
    + \int_0^1 b(x)\, u'(x)\, \varphi(x)\dd x
    + \int_0^1 c(x)\, u(x)\, \varphi(x)\dd x
    = \int_0^1 f(x)\, \varphi(x)\dd x.
\end{equation}

\paragraph{Query tool (\texttt{QUERY}).}
The model submits a test function $\varphi(x)$ as a symbolic expression
and receives the scalar $\int_0^1 f(x)\,\varphi(x)\dd x$.  Each query
costs one unit from a finite budget (50 for the hard problem).

\paragraph{Weight integral tool (\texttt{COMPUTE: term\_integral}).}
For any test function $\varphi$, the model can request the three weight
integrals from the left-hand side of~\eqref{eq:weak}:
\begin{align*}
    \text{diffusion:} &\quad \textstyle\int_0^1 u'(x)\,\varphi'(x)\dd x, \\
    \text{advection:} &\quad \textstyle\int_0^1 u'(x)\,\varphi(x)\dd x, \\
    \text{reaction:}  &\quad \textstyle\int_0^1 u(x)\,\varphi(x)\dd x.
\end{align*}
These are computable from the known solution $u(x)$ and the test function
$\varphi(x)$ alone, with no knowledge of $a$, $b$, $c$.  This tool is free
and unlimited.

\paragraph{Solve tool (\texttt{COMPUTE: solve}).}
After accumulating at least $4N$ queries, the model can request an
automatic least-squares recovery of all coefficients from the stored query
data.  This is free.

\paragraph{Verify tool (\texttt{COMPUTE: verify}).}
The model can test its current coefficient estimates against a new test
function by comparing the predicted left-hand side of~\eqref{eq:weak}
(computed from estimated $\hat{a}, \hat{b}, \hat{c}$) to the actual
$\int f\varphi\dd x$ (which costs one query).

\subsection{What the Model Must Figure Out}

The model must decide:
\begin{enumerate}[nosep]
    \item Which test functions to submit as queries.
    \item When to call \texttt{solve} and whether to add more queries
          afterward.
    \item How to use \texttt{term\_integral} to evaluate whether a
          candidate test function is informative \emph{before} spending
          a query on it.
    \item When the coefficients are accurate enough to submit.
\end{enumerate}

\noindent The critical scientific insight is that not all test functions
contribute equally to the linear system.

%----------------------------------------------------------------------
\section{The Ill-Conditioning Problem}
%----------------------------------------------------------------------

The weak form~\eqref{eq:weak} produces one linear equation per test
function, relating the $4N$ unknown coefficients to the measured
$\int f\varphi\dd x$.  Each equation has three contributions on the
left-hand side---diffusion, advection, and reaction---weighted by the
integrals that the \texttt{term\_integral} tool returns.

For a typical test function, these weights have very different magnitudes.
The diffusion weight involves $u'(x)\,\varphi'(x)$, where both factors
are $O(1)$ for smooth functions.  The reaction weight involves
$u(x)\,\varphi(x)$, and since the solution $u(x)$ satisfies homogeneous
boundary conditions and is often small in magnitude ($|u| \sim 0.01$--$0.1$
in our problems), this weight is correspondingly small.  As a result:
\begin{itemize}[nosep]
    \item The diffusion coefficient $a(x)$ dominates every equation and is
          easy to recover.
    \item The advection coefficient $b(x)$ has medium influence and is
          recovered with moderate accuracy.
    \item The reaction coefficient $c(x)$ has tiny influence and is
          difficult to recover---errors in $a$ and $b$ are absorbed into
          the estimate of $c$.
\end{itemize}
This is a well-known property of elliptic inverse problems: the reaction
term is the hardest to identify from weak-form data.  A mathematician
would compensate by choosing test functions that amplify the reaction
weight relative to the diffusion weight---for example, very smooth
functions with small derivatives but large values where $|u(x)|$ is
maximized.

The benchmark tests whether the model can discover and act on this insight.

%----------------------------------------------------------------------
\section{Experimental Design}
%----------------------------------------------------------------------

\subsection{The Controlled Variable}

We compare two conditions that differ \textbf{only in the initial user
message} sent to the model.  All code, tools, PDE instances, random seeds,
and query budgets are identical.

\paragraph{Baseline prompt.}
\begin{quote}
\texttt{Begin.}
\end{quote}

\paragraph{Enhanced prompt.}
\begin{quote}
\texttt{Begin. Important background: the weak form equation is}

\smallskip
\texttt{~~$\int a(x)u'(x)\varphi'(x)\,dx + \int b(x)u'(x)\varphi(x)\,dx
  + \int c(x)u(x)\varphi(x)\,dx = \int f(x)\varphi(x)\,dx$}

\smallskip
\texttt{The first term (diffusion) involves $u'$ and $\varphi'$. The second
  term (advection) involves $u'$ and $\varphi$. The third term (reaction)
  involves $u$ and $\varphi$.}

\smallskip
\texttt{Each test function $\varphi$ produces one equation with three unknown
  contributions on the left. If one contribution dominates (say the
  diffusion weight is 100x larger than the reaction weight), then that
  equation primarily constrains the diffusion coefficient $a(x)$ and tells
  you almost nothing about $c(x)$. An ill-conditioned system arises when
  all your equations are dominated by the same term.}

\smallskip
\texttt{Use COMPUTE: term\_integral to check the three weights for a test
  function before spending a query on it. Look for test functions that give
  balanced weights across the three terms, or that emphasize a term that
  your existing equations have not yet constrained well.}
\end{quote}

\noindent The enhanced prompt explains \emph{why} the system can be
ill-conditioned and \emph{which tool} to use to diagnose it, but does not
specify which test functions to choose.  The model must discover those on
its own.

\subsection{Evaluation Metrics}

We report:
\begin{itemize}[nosep]
    \item \textbf{Coefficient error}: $\max_j |c_j^{\text{pred}} -
          c_j^{\text{true}}|$ for each coefficient function, and total.
    \item \textbf{Pointwise error}: $\max_{x \in [0,1]}
          |c^{\text{pred}}(x) - c^{\text{true}}(x)|$ for each coefficient.
    \item \textbf{Queries used}: total number of budget-consuming actions.
    \item \textbf{Efficiency AUC}: area under the error-vs-queries curve,
          computed post hoc by solving the system at each query count.
\end{itemize}

%----------------------------------------------------------------------
\section{Results}
%----------------------------------------------------------------------

\subsection{Hard Difficulty, Seed 42}

Table~\ref{tab:results} compares the baseline and enhanced prompts on the
same hard problem instance ($N = 4$, 16 unknowns).

\begin{table}[h]
\centering
\caption{Coefficient recovery on hard difficulty, seed 42.}\label{tab:results}
\begin{tabular}{lcc}
\toprule
\textbf{Metric} & \textbf{Baseline} & \textbf{Enhanced} \\
\midrule
Total coefficient error & 1.78 & \textbf{1.03} \\
Total pointwise error   & 2.79 & \textbf{1.58} \\
$a(x)$ coeff error      & 0.051 & 0.021 \\
$b(x)$ coeff error      & 0.364 & 0.182 \\
$c(x)$ coeff error      & 1.365 & 0.824 \\
$f(x)$ coeff error      & 0.000 & 0.000 \\
Queries used            & 17   & 17 \\
Efficiency AUC          & 1.96 & \textbf{1.03} \\
Queries to error $< 1.0$ & never & 8 \\
Test function diversity & 9 poly, 7 trig & 7 poly, 8 trig, 1 localized \\
\bottomrule
\end{tabular}
\end{table}

The enhanced prompt reduces total coefficient error by 42\% and pointwise
error by 43\%.  The improvement is concentrated in the hardest-to-recover
coefficients: $c(x)$ error drops by 40\% and $b(x)$ error by 50\%.
Both runs recover $f(x)$ perfectly and $a(x)$ to high accuracy.

\subsection{Qualitative Differences in Reasoning}

The most striking difference is not the final error but the \emph{reasoning
process}.  The baseline model selects test functions mechanically, never
checking whether they are informative:

\begin{quote}
\emph{Baseline, Turn 18:} ``Let me try another polynomial test function.''
$\to$ queries $x(1{-}x)^2$ without checking weights.
\end{quote}

The enhanced model checks weights before every query and articulates the
balance:

\begin{quote}
\emph{Enhanced, Turn 11:} ``The diffusion still dominates but only by about
6x, which is better.''

\emph{Enhanced, Turn 36:} ``The advection and reaction terms are almost
equal in magnitude, and both are only about 5x smaller than
diffusion\ldots\ This should give me good information about both advection
and reaction coefficients.''

\emph{Enhanced, Turn 80:} ``The advection and reaction terms are almost
exactly equal, and diffusion is only about 2.6x larger. This is very well
balanced.''
\end{quote}

The enhanced model found test functions like $x^2(1{-}x)^3$, where the
three weights are $-0.00424$, $-0.00161$, $-0.00161$---genuinely balanced.
The baseline never discovered such functions because it never looked for
them.

%----------------------------------------------------------------------
\section{Discussion}
%----------------------------------------------------------------------

\subsection{What the Benchmark Measures}

This benchmark occupies a specific point in the space of LLM evaluations:

\begin{enumerate}[nosep]
    \item \textbf{Not knowledge recall.} The model is not asked to recite
          facts about PDEs.  It must \emph{use} the weak form to design
          experiments.
    \item \textbf{Not tool use alone.} Both conditions have identical
          tools.  The difference is whether the model uses them
          \emph{strategically}.
    \item \textbf{Not arithmetic.} The model never solves a linear system
          itself; \texttt{COMPUTE: solve} handles that.  The model's job
          is experimental design.
    \item \textbf{Scientific metacognition.} The model must reason about
          \emph{which measurements are informative} and why, which is the
          core skill of experimental science.
\end{enumerate}

\subsection{The Ill-Conditioning Barrier}

The reaction coefficient $c(x)$ remains the hardest to recover in both
conditions.  This is not a limitation of the benchmark but a genuine
property of elliptic inverse problems: the reaction term contributes
negligibly to the weak form when $|u(x)|$ is small.  Even a perfect
experimental design cannot fully overcome this when the solution magnitude
is small relative to its derivative.

The enhanced prompt reduces $c(x)$ error by 40\% but does not eliminate it.
This suggests that the remaining error is dominated by the inherent
ill-conditioning rather than by suboptimal test function selection.

\subsection{One Paragraph, 42\% Improvement}

The enhanced prompt is a single paragraph of mathematical context.  It does
not tell the model which test functions to use, does not provide a
strategy, and does not mention specific functions like Gaussians or
high-order polynomials.  It explains a \emph{principle}---that term
dominance causes ill-conditioning---and points to a \emph{tool} for
diagnosing it.

The model translates this principle into action: it checks weights before
querying, seeks balanced functions, and articulates its reasoning in
mathematical terms.  This is the kind of capability transfer that
characterizes genuine understanding versus pattern matching.

%----------------------------------------------------------------------
\section{Conclusion}
%----------------------------------------------------------------------

We have presented a benchmark that tests scientific reasoning in language
models through the concrete task of PDE identification.  The benchmark
provides the model with the tools of a mathematician---weak-form
integrals, weight decomposition, least-squares solving---and measures
whether it can assemble them into an effective experimental strategy.

Our main finding is that a single paragraph of mathematical context about
ill-conditioning produces a 42\% improvement in coefficient recovery, with
the model exhibiting qualitatively different reasoning about experimental
design.  This demonstrates that current language models can engage in
meaningful scientific reasoning when given appropriate conceptual
scaffolding, but cannot independently discover the key insight that makes
the problem tractable.

The gap between the baseline and enhanced conditions---both using identical
tools on the same problem---provides a clean, reproducible measure of
scientific metacognition in language models.

\end{document}
